% !TEX program = lualatex
% Auto-generated from YAML (v3) — 09: データサイエンス・AI（2）：活用事例と倫理的課題

\documentclass[aspectratio=43,professionalfonts,handout,10pt]{beamer}
\usepackage{beamer_template}

\newcommand{\LectureNum}{09}
\newcommand{\LectureTitle}{データサイエンス・AI（2）：活用事例と倫理的課題}
\newcommand{\CourseName}{情報リテラシー}
\newcommand{\TextPages}{pp.124-127}
\newcommand{\NextTopic}{データサイエンス・AI（3）：データ分析と仮説検証サイクル}
\newcommand{\NextPages}{pp.126-129}

\begin{document}
% --- Slide 1: title ---

\begin{frame}{\CourseName\quad 第\LectureNum 回「\LectureTitle 」}
  {\small \TermName\quad \TeacherName\quad ／\quad 教科書 \TextPages}

  \vfill

  \begin{block}{今日の流れ}
    \begin{enumerate}
      \item データの種類（質的・量的）と尺度水準
      \item データの可視化・テキストマイニング
      \item 事例1：採用AIとバイアス
      \item 事例2：顔認識とプライバシー
      \item 事例3：ブラックボックス問題
    \end{enumerate}
  \end{block}
  \vfill

  \begin{block}{今日の目標}
    \begin{itemize}
      \item データサイエンスやAI技術が課題解決に役立ち，専門知識と組み合わせて価値を生み出すことを事例を通じて説明できる。
    \end{itemize}
  \end{block}
  \vfill

  \begin{exampleblock}{キーワード}
    \small \textbf{欠損値} ／ \textbf{外れ値・異常値} ／ \textbf{データクレンジング}
  \end{exampleblock}
\end{frame}

% --- Slide 2: twocol ---

\begin{frame}{データの2つの種類}
  \begin{columns}[T]
    \begin{column}{0.48\textwidth}
      \begin{block}{質的データ}
        \small
        \begin{itemize}
          \item テストの順位や血液型など，種類を区別するデータ
          \item 数量で表すことができない
          \item 例：性別，血液型，出身地，通学手段
        \end{itemize}
      \end{block}
    \end{column}
    \begin{column}{0.48\textwidth}
      \begin{exampleblock}{量的データ}
        \small
        \begin{itemize}
          \item 気温や体重など，数値で表すことができるデータ
          \item 大小比較や計算が意味を持つ
          \item 例：身長，体重，気温，テストの点数，通学時間
        \end{itemize}
      \end{exampleblock}
    \end{column}
  \end{columns}
\end{frame}

% --- Slide 3: table ---

\begin{frame}{尺度水準（データの種類）}
  \centering
  \begin{tabular}{llll}
    \toprule
    尺度 & 特徴 & 例 & できる演算 \\
    \midrule
    名義尺度 & 数値の順序や大小に意味はない & 通学手段・血液型 & 度数・最頻値 \\
    順序尺度 & 大小には意味があるが，間隔が一定でない & 好きな行事の順位・満足度 & 中央値・最頻値 \\
    間隔尺度 & 等間隔だが，比には意味がない（ゼロ点が定まらない） & 西暦・気温（℃）・偏差値 & 平均値・標準偏差 \\
    比例尺度 & 等間隔かつ絶対的なゼロ点がある & 身長・体重・通学時間・金額 & すべて可能 \\
    \bottomrule
  \end{tabular}
\end{frame}

% --- Slide 5: points ---

\begin{frame}{データの可視化とテキストマイニング}
  \begin{block}{ポイント}
    \begin{itemize}
      \item データは，集めただけでは数値や文字情報でしかない。分析・関連付けで初めて情報としての価値が現れる
      \item 可視化：表形式で整理→適切なグラフで表すことで全体像が把握できる
      \item テキストマイニング：文章を構成する単語などに着目し，出現回数・傾向・相関などを解析する
      \item ワードクラウド：テキストデータ中の単語の出現頻度を視覚化（頻出単語が大きく表示される）
      \item 活用例：SNSの商品イメージ分析，自由記述アンケートの傾向分析，AIの自然言語処理の基盤技術
      \item 形態素解析：文章を最小意味単位（形態素）に分割して品詞を判定する処理。テキストマイニングの基盤技術
    \end{itemize}
  \end{block}
\end{frame}

% --- Slide 6: points ---

\begin{frame}{事例1：採用AIとバイアス}
  \begin{block}{ポイント}
    \begin{itemize}
      \item ある大手IT企業が採用AIを開発 → 女性を低評価してしまう問題が発覚 → AIの使用を中止
      \item なぜ起きたか：過去10年のデータが男性中心 → AIが「男性＝良い候補者」と誤学習
      \item メリット：効率化（数千件を短時間で処理）・客観的評価・コスト削減
      \item 重要：データの偏り ＝ AIの偏り。AIは過去の偏見を未来に持ち込む
    \end{itemize}
  \end{block}

  \vfill

  \begin{exampleblock}{補足}
    \begin{itemize}
      \item 今日のColab実習：年齢層別の採用率を分析し，バイアスを体験する
    \end{itemize}
  \end{exampleblock}
\end{frame}

% --- Slide 7: twocol ---

\begin{frame}{事例2：顔認識とプライバシー}
  \begin{columns}[T]
    \begin{column}{0.48\textwidth}
      \begin{block}{活用例}
        \small
        \begin{itemize}
          \item 空港の入国審査
          \item スマホのロック解除（Face ID）
          \item 行方不明者の捜索
          \item 商業施設の防犯
        \end{itemize}
      \end{block}
    \end{column}
    \begin{column}{0.48\textwidth}
      \begin{exampleblock}{リスクとプライバシー問題}
        \small
        \begin{itemize}
          \item 常に追跡される不安・行動の自由が制限される
          \item 個人情報保護法（同意が必要）
          \item 誤認識のリスク（冤罪の可能性）
          \item データ漏洩のリスク
        \end{itemize}
      \end{exampleblock}
    \end{column}
  \end{columns}
\end{frame}

% --- Slide 8: points ---

\begin{frame}{事例3：ブラックボックス問題}
  \begin{block}{ポイント}
    \begin{itemize}
      \item 説明可能性：AIがなぜその判断をしたかが分からない。複雑なモデルほど説明が難しい
      \item なぜ問題か：判断の理由が説明できないと信頼されない／間違いに気づけない／責任の所在が不明確
      \item 問題となる場面：医療診断AI（なぜその診断？）／融資審査AI（なぜ却下？）／採用判定AI（なぜ不採用？）
      \item 対策：説明可能なAI（XAI: Explainable AI）の研究が進んでいる
    \end{itemize}
  \end{block}
\end{frame}

% --- Slide 9: practice ---

\begin{frame}{Colab実習：採用データの年齢層別分析}
  {\small 実際にコンピュータで操作してください。}

  \vfill

  \begin{block}{手順}
    \begin{enumerate}
      \item TeamsフォルダリンクをクリックしてColabを開く
      \item 「採用データ.csv」を自分のマイドライブにコピー
      \item セル1を実行（Googleドライブへのアクセス許可）
      \item Part 2のセル3とセル4を完成させる
      \item グラフが表示されたらスクリーンショット（Win + Shift + S）
      \item Teamsの課題に提出
    \end{enumerate}
  \end{block}
\end{frame}

% --- チェックリスト ---

\begin{frame}{まとめ・チェックリスト}
  \begin{block}{自己チェック（できたら $\checkmark$ をつけよう）}
    \begin{itemize}
      \item[$\square$] 採用AI・顔認識など実際の活用事例からAIのバイアス問題を学んだ
      \item[$\square$] データの可視化とテキストマイニングの基本を理解した
      \item[$\square$] ブラックボックス問題とAI倫理の重要性を考えた
    \end{itemize}
  \end{block}

  \vfill

  {\small 質問があれば授業後またはTeamsで受け付けます。}
\end{frame}

\end{document}
